\chapter{Seguridad y Row Level Security}\label{cap:seguridad}

\section{Modelo de amenazas}

Se consideran las amenazas más relevantes del OWASP Top 10:

\begin{itemize}[leftmargin=*,itemsep=0.2em]
  \item Acceso no autorizado a recursos de otros usuarios
        (\emph{broken access control}).
  \item Inyección SQL (\emph{injection}).
  \item Exposición de datos sensibles (\emph{cryptographic failures}).
  \item \emph{Cross-Site Scripting} (XSS).
  \item Configuración insegura (\emph{security misconfiguration}).
  \item \emph{Server-Side Request Forgery} en el agente.
\end{itemize}

\section{Autenticación y autorización}

La autenticación se delega en Supabase Auth (correo/contraseña +
Google OAuth). Cada petición autenticada al backend incluye un JWT en
\texttt{Authorization: Bearer <token>}, verificado contra las claves
públicas del proyecto Supabase. El \emph{claim} \texttt{sub} del JWT
identifica al usuario y se propaga a la consulta SQL como
\texttt{auth.uid()}.

\section{Row Level Security}

Todas las tablas con datos de usuario activan RLS. La política básica
es:

\begin{lstlisting}[language=SQL,caption={Politica RLS por usuario.}]
CREATE POLICY user_isolation
  ON public.analyses
  FOR ALL
  USING (auth.uid() = user_id)
  WITH CHECK (auth.uid() = user_id);
\end{lstlisting}

Las consultas administrativas usan la \emph{service role key} y
\texttt{set local role = service\_role}.

\section{Protección frente a XSS y abuso del DOM}

\begin{itemize}[leftmargin=*,itemsep=0.2em]
  \item En React, todo el texto dinámico se renderiza con
        \texttt{\{variable\}}, que escapa por defecto.
  \item En la extensión, la composición del informe de verificación
        usa \texttt{textContent} y \texttt{createElement}, evitando
        \texttt{innerHTML}.
  \item Los enlaces externos llevan \texttt{rel="noopener noreferrer"}
        y \texttt{target="\_blank"}.
\end{itemize}

\section{Inyección y validación}

La validación tipada de Pydantic en el backend rechaza payloads
malformados. Las consultas a Supabase usan PostgREST con parámetros
nombrados; no se construye SQL con concatenación de cadenas.

\section{Gestión de secretos}

Todos los secretos (\texttt{SUPABASE\_SERVICE\_ROLE\_KEY},
\texttt{STRIPE\_SECRET\_KEY}, \texttt{STRIPE\_WEBHOOK\_SECRET},
\texttt{TAVILY\_API\_KEY}) se inyectan vía variables de entorno y se
documentan en \path{.env.example}. El repositorio nunca incluye
secretos reales.

\section{RGPD: borrado de cuenta}

El endpoint \texttt{POST /account/delete} elimina en cascada los datos
del usuario y solicita la baja en Supabase Auth, cumpliendo el
derecho al olvido.
